% Fate's Edge Quickstart
% Compile with pdflatex
\documentclass[10pt, twocolumn]{article}
\usepackage[utf8]{inputenc}
\usepackage[T1]{fontenc}
\usepackage{geometry}
\geometry{margin=1in, columnsep=0.5in}
\usepackage{graphicx}
\usepackage{parskip}
\usepackage{tabularx}
\usepackage{booktabs}
\usepackage{xcolor}
\usepackage{titlesec}
\titleformat{\section}{\large\bfseries}{\thesection}{1em}{}
\titleformat{\subsection}{\bfseries}{\thesubsection}{1em}{}
\setlength{\parindent}{0pt}
\setlength{\parskip}{0.8ex plus 0.3ex minus 0.2ex}
\pagestyle{plain}

% Simple box for examples
\newenvironment{example}{\color{black!70!green}\small\textbf{Example: }}{\par}

\begin{document}

\twocolumn[
\centering
{\huge\bfseries Fate's Edge Quickstart}\\[0.5cm]
{\large Your first session in 15 minutes}\\[0.3cm]
\rule{\linewidth}{0.5pt}\\[0.5cm]
]

\section*{Welcome to Fate's Edge}
You are exiles, rogues, and survivors in a world where every choice carries weight and every roll of the dice can change the story. This Quickstart contains everything you need to play your first session:
\begin{itemize}
\item Core rules (2 pages)
\item 4 pre‑generated characters
\item A complete starter adventure: \textit{The Lantern at Dusk}
\item Game Master tips
\end{itemize}
No other books required.

\section{Core Resolution}
When your character attempts something risky or uncertain:
\begin{enumerate}
\item \textbf{Declare} what you want to achieve and how (\textit{Attribute + Skill}).
\item \textbf{GM sets Difficulty Value (DV)}: 2 (easy), 3 (standard), 4 (hard), 5+ (extreme).
\item \textbf{GM sets Position}: Dominant (safe), Controlled (normal), Desperate (dangerous).
\item \textbf{Roll} a number of d10s = Attribute + Skill.
\item \textbf{Count successes} (each 6--10 = 1 success). Each 1 gives the GM a \textbf{Story Beat (SB)}.
\item \textbf{Compare successes to DV} using the table below.
\end{enumerate}

\subsection*{Outcome Matrix}
\begin{tabularx}{\linewidth}{X c X}
\toprule
Result & Condition & Effect \\
\midrule
Clean Success & successes $\ge$ DV, SB=0 & You do it, no complication. \\
Success with SB & successes $\ge$ DV, SB$\ge$1 & You do it, GM spends SB for a twist. \\
Partial & $0<$ successes $<$ DV & You make progress, gain 1 Boon. \\
Miss & successes $=$ 0 & You fail, gain 2 Boons. \\
\bottomrule
\end{tabularx}

\subsection*{Boons}
Boons represent luck and momentum. Hold up to 5. At the end of a scene, reduce to 2 (lost excess). Spend 1 Boon to:
\begin{itemize}
\item Re‑roll a single die after seeing the pool.
\item Improve your Position by one step before rolling.
\item Activate an Asset (see character sheets).
\item Convert 2 Boons $\to$ 1 XP (once per session, max 2 XP).
\end{itemize}

\subsection*{Story Beats (SB) -- GM Only}
Spend SB to introduce complications:
\begin{itemize}
\item 1 SB -- minor: noise, dropped item, tick a timer.
\item 2 SB -- moderate: alarm raised, lose cover.
\item 3+ SB -- serious: reinforcements, scene shift.
\end{itemize}

\section{Combat in Brief}
Same as above, with these additions:
\begin{itemize}
\item \textbf{Turn order}: Narrative initiative -- GM spotlights based on fiction. Ambushers act first.
\item \textbf{Move}: Shift one range band (Close $\leftrightarrow$ Near $\leftrightarrow$ Far) as a free movement.
\item \textbf{Action}: Attack, cast a spell, use a talent, etc.
\item \textbf{Opportunity attack}: Leaving Close range without using \textit{Disengage} provokes one free melee attack.
\item \textbf{Fatigue}: Each point of Fatigue worsens Position one step (Dominant$\to$Controlled$\to$Desperate). If already Desperate, each Fatigue gives --1 die. Full rest removes all.
\item \textbf{Harm}: Harm 1 (--1 die), Harm 2 (--2 dice), Harm 3 (incapacitated).
\end{itemize}

\section{Attributes \& Skills}
\begin{tabularx}{\linewidth}{XXX}
\toprule
Attribute & What it does & Example Skills \\
\midrule
Body & strength, agility, endurance & Melee, Athletics, Stealth \\
Wits & perception, quick thinking & Insight, Lore, Investigation \\
Spirit & willpower, intuition & Resolve, Arcana, Survival \\
Presence & charisma, leadership & Sway, Command, Deception \\
\bottomrule
\end{tabularx}
\textbf{Skill ratings}: 0 untrained, 1 novice, 2 competent, 3 expert, 4 master, 5 legendary.

\section{Pre‑Generated Characters}
Choose one. All are balanced for the starter adventure.

\subsection*{1. Kael -- Soldier}
\textit{Once a legionary, now a sellsword.}\\
\textbf{Body 3, Wits 2, Spirit 2, Presence 1}\\
\textbf{Skills}: Melee 3, Athletics 2, Endurance 1, Command 1\\
\textbf{Talent -- Battle Instincts}: Once per scene, reroll a failed defense roll.\\
\textbf{Asset -- Trusted Blade}: Your sword never breaks on a Miss.\\
\textbf{Bonds}: (choose one with another PC) "I saved your life outside Millhaven."

\subsection*{2. Lyra -- Scholar}
\textit{Expelled from the Academy for forbidden research.}\\
\textbf{Wits 3, Spirit 3, Body 1, Presence 1}\\
\textbf{Skills}: Lore 2, Arcana 2, Investigation 2, Insight 1\\
\textbf{Talent -- Lorekeeper}: Once per session, recall obscure history without rolling.\\
\textbf{Asset -- Scholar's Satchel}: Contains books, reagents, a small telescope.\\
\textbf{Bonds}: "We both lost our place in Silkstrand."

\subsection*{3. Sera -- Shadow}
\textit{Former guild thief, now a wanderer.}\\
\textbf{Wits 3, Body 2, Spirit 2, Presence 1}\\
\textbf{Skills}: Stealth 3, Subterfuge 2, Streetwise 1, Melee 1\\
\textbf{Talent -- Silver Tongue}: +1 die to persuade or deceive through speech.\\
\textbf{Asset -- Shadow's Cloak}: +1 die to stealth rolls in dim light.\\
\textbf{Bonds}: "You kept my secret when the guards came."

\subsection*{4. Mira -- Healer}
\textit{Accused of witchcraft for curing the incurable.}\\
\textbf{Spirit 3, Wits 2, Body 2, Presence 1}\\
\textbf{Skills}: Medicine 3, Survival 2, Insight 2, Nature 1\\
\textbf{Talent -- Iron Stomach}: Immune to mundane poisons and spoiled food.\\
\textbf{Asset -- Healer's Kit}: Bandages, herbs, basic medical supplies.\\
\textbf{Bonds}: "You trusted me when no one else would."
\vspace{0.2cm}

\noindent Each PC starts with \textbf{3 Boons} and \textbf{0 Fatigue/Harm}.

\section{Starter Adventure -- The Lantern at Dusk}
\subsection*{Setup}
The party is hired by Elder Sarai of the remote village of \textbf{Duskwood}. A terrible blight is spreading from the old barrow in the hills -- the same barrow where a cursed lantern was buried centuries ago. The lantern must be brought to a stone circle in the valley to break the curse. But the barrow is not empty.

\subsection*{What the Party Knows}
\begin{itemize}
\item The lantern is in the central chamber of the barrow.
\item The barrow is guarded by restless spirits (the \textbf{Duskwights}).
\item A rival treasure hunter, \textbf{Quick Lena}, has also entered the barrow.
\item The curse will kill the village's crops in three days.
\end{itemize}

\subsection*{Timers}
\begin{itemize}
\item \textbf{Barrow Collapse [6]} -- ticks when the party fails a roll inside (GM advances 1 segment per Miss). At 6, the entrance seals -- escape becomes a Desperate group action.
\item \textbf{Lena's Agenda [4]} -- ticks at the end of each scene inside the barrow. When it fills, Lena has already taken the lantern; the party must chase her.
\end{itemize}

\subsection*{Scene 1 -- The Entry}
The barrow entrance is a narrow stone passage. A faint green glow pulses from deeper inside.  
\textbf{Obstacle}: A cave‑in blocks the way (Body+Athletics to climb, DV 3, Controlled).  
\textit{Partial}: You get through but disturb rubble -- tick Barrow Collapse +1.  
\textit{Miss}: You are pinned by a rock -- mark Harm 1 and tick Collapse +1.

\subsection*{Scene 2 -- The Spirit Corridor}
The passage widens into a hall where three Duskwights drift. They do not attack unless provoked, but they will not let anyone pass without an offering.  
\textbf{Offer a memory or a secret} (Presence+Sway, DV 3, Controlled).  
\textit{Clean Success}: They bow and let you pass.  
\textit{Partial}: They take your offering but also steal a random piece of equipment (GM chooses).  
\textit{Miss}: They attack as a single mob (Cap 2, Harm 1 each, 2 Fatigue each on hit). Combat follows normal rules.

\subsection*{Scene 3 -- The Lantern Chamber}
The lantern sits on a stone altar. Quick Lena is there, trying to pry it loose.  
\textbf{Confront Lena}: She is not a fighter. Offer a deal or intimidate her (Presence+Sway, DV 3, Controlled).  
\textit{Success}: She leaves peacefully, but warns that "others will come for it."  
\textit{Partial}: She agrees to split any reward but demands a favor later (start a Debt Timer [4]).  
\textit{Miss}: She triggers a trap -- a stone door starts closing (tick Collapse +2).

\subsection*{Scene 4 -- Escape}
After you take the lantern, the barrow begins to collapse in earnest.  
\textbf{Group skill challenge}: The party needs 3 successes before 2 failures to escape. Each PC describes how they help (e.g., Athletics to brace a door, Lore to read ancient escape runes). DV 3, Position Desperate.  
Each failure ticks Collapse +1. If Collapse reaches 6 before you escape, the entrance seals -- you must find another way (start a second Timer [4] to dig out, each success removes 1 segment; on failure, mark Fatigue).  
\textit{Success}: You burst into the moonlight, lantern in hand.  
\textit{Failure}: You escape, but someone is left behind -- captured or lost. The remaining party must rescue them later.

\subsection*{Resolution}
Return the lantern to the stone circle. Breaking the curse requires a simple ritual (Spirit+Lore, DV 2).  
\textit{Success}: The blight lifts. Elder Sarai rewards each PC with a \textbf{String} -- "The village owes you."  
\textit{Partial}: The curse lifts, but the lantern cannot be destroyed -- it will attract another danger soon.  
\textit{Miss}: The curse is only pushed back one year. A dark spirit escapes and will become a recurring enemy.

\subsection*{Advancement}
Each PC earns \textbf{6--10 XP} at the end. XP can be spent:
\begin{itemize}
\item Raise an Attribute: new rating $\times$ 3 XP (downtime = new rating days)
\item Raise a Skill: new level $\times$ 2 XP (downtime = new level days)
\item Learn a Minor Talent (2--4 XP) -- examples: \textit{Keen Senses}, \textit{Second Wind}, \textit{Combat Reflexes}
\item Acquire a Minor Asset (4 XP) -- safehouse, contact, or special tool.
\end{itemize}
\textbf{Rush rule}: Skip downtime by accepting a \textit{Haste timer [4]}. If it fills before you finish, the new ability has a flaw (GM's choice).

\section{GM Tips -- Running Your First Game}
\begin{itemize}
\item \textbf{Start strong}: Use the cold open -- "The barrow's green light flickers. You smell rot. What do you do?"
\item \textbf{Set DV to 3} by default. Only raise to 4 or 5 if the task is truly legendary.
\item \textbf{Spend SB early and often}. A single SB can turn a quiet scene into a tense one.
\item \textbf{Fail forward}: On a Miss, never say "nothing happens." Always introduce a new complication or advance a timer.
\item \textbf{Use the pre‑gens} -- they are designed to showcase different play styles.
\item \textbf{Keep the spotlight moving}. Ask quiet players directly: "Sera, what are you doing while this happens?"
\item \textbf{End on a cliffhanger}. When the lantern is lifted, describe the ceiling cracking -- then say "to be continued."
\end{itemize}

\section*{Optional -- Magic in One Page}
If a player wants a magical character, use these simplified rules (based on the \textit{Free Caster} path).

\subsection*{Spellcasting}
Choose a magical \textbf{Element} (Fire, Water, Earth, Air, Fate, Life, Luck, Death). Describe your spell using \textbf{TAGS}. Each tag increases DV by 1. Dangerous tags (Teleport, Transform, Dominate) increase DV by +2 instead.  
\textbf{DV = 1 + number of tags} (max 6). Roll \textbf{Spirit + Arcana}.

Example: \textit{"I throw a [Burning] [Area] blast"} = DV 3.

\subsection*{Backlash}
If you roll a Miss or the GM spends SB, magic backfires:
\begin{itemize}
\item 1--2 SB: minor flare, Fatigue 1, eerie sound.
\item 3+ SB: Harm 1, the spell affects an ally, or a Patron omen appears.
\end{itemize}

\subsection*{Sample Spells (DV 2--4)}
\begin{itemize}
\item \textit{Light} [Illuminate] -- DV 2. Create a small, sustained light.
\item \textit{HEAL} [HEAL] -- DV 2. Heal 1 Fatigue from yourself or a touched ally.
\item \textit{Shield of Air} [Ward] [Defend] -- DV 3. Gain Position +1 on next defense.
\item \textit{Flame Cloak} [Burning] [Area] -- DV 3. Enemies suffer --1 die to approach.
\end{itemize}
For full magic systems (Runekeepers, Invokers, Cantors), see the complete \textit{Fate's Edge Player's Guide}.

\section*{License \& Credits}
\textit{Fate's Edge Quickstart} \textcopyright{} Nicholas A. Gasper. This work is licensed under the Creative Commons Attribution 4.0 International License. \texttt{http://creativecommons.org/licenses/by/4.0/}  
Based on the \textit{Fate's Edge} roleplaying game.  
Cover illustration placeholder.

\section*{Ready for More?}
This Quickstart gives you a complete first session. For the full experience:
\begin{itemize}
\item \textbf{Player's Guide}: Attributes, Skills, Talents, magic systems, character advancement, and the world.
\item \textbf{GM Guide}: Timers, fronts, advanced Story Beat management, faction rules, and high‑tier play.
\item \textbf{Adventures}: \textit{Blood and Silk} (village defense), \textit{Beyond Millhaven} (mediation), \textit{The Brass Gate} (heist).
\end{itemize}
All available under CC‑BY from the author. Now gather your dice and step into the Edge.

\end{document}
